% =============================================================================
% bab4-implementasi.tex — BAB IV IMPLEMENTASI PROYEK BISNIS/STUDI KASUS
% Proyek Bisnis / Studi Kasus — Lampiran II (S2)
% =============================================================================

\chapter{IMPLEMENTASI PROYEK BISNIS/STUDI KASUS}
\label{bab:implementasi}

\section{Analisis Kondisi Bisnis Saat Ini}
\label{sec:analisis-kondisi}

Uraikan hasil analisis kondisi bisnis saat ini berdasarkan data yang telah dikumpulkan.
Gunakan alat analisis yang telah ditentukan pada Bab III untuk mendeskripsikan kondisi
[aspek bisnis] pada [nama organisasi].

\subsection{Analisis Faktor Internal}
\label{subsec:analisis-internal}

Uraikan hasil analisis faktor internal organisasi (kekuatan dan kelemahan):

\begin{table}[H]
\caption{Matriks Analisis Faktor Internal}
\label{tab:analisis-internal}
\begin{isitable}
\begin{tabular}{p{0.5cm} p{5cm} p{5cm} p{4cm}}
\toprule
\textbf{No} & \textbf{Kekuatan (Strengths)} & \textbf{Kelemahan (Weaknesses)} & \textbf{Implikasi} \\
\midrule
1 & Kekuatan pertama & Kelemahan pertama & Implikasi \\
2 & Kekuatan kedua & Kelemahan kedua & Implikasi \\
\bottomrule
\end{tabular}
\end{isitable}
\end{table}

\subsection{Analisis Faktor Eksternal}
\label{subsec:analisis-eksternal}

Uraikan hasil analisis faktor eksternal (peluang dan ancaman):

\begin{table}[H]
\caption{Matriks Analisis Faktor Eksternal}
\label{tab:analisis-eksternal}
\begin{isitable}
\begin{tabular}{p{0.5cm} p{5cm} p{5cm} p{4cm}}
\toprule
\textbf{No} & \textbf{Peluang (Opportunities)} & \textbf{Ancaman (Threats)} & \textbf{Implikasi} \\
\midrule
1 & Peluang pertama & Ancaman pertama & Implikasi \\
2 & Peluang kedua & Ancaman kedua & Implikasi \\
\bottomrule
\end{tabular}
\end{isitable}
\end{table}

\section{Identifikasi dan Analisis Permasalahan Bisnis}
\label{sec:identifikasi-masalah-bisnis}

Berdasarkan hasil analisis kondisi bisnis saat ini, permasalahan utama yang diidentifikasi
adalah sebagai berikut:

\begin{enumerate}
  \item \textbf{Permasalahan Pertama:} Deskripsi permasalahan dan dampaknya terhadap
        kinerja bisnis.
  \item \textbf{Permasalahan Kedua:} Deskripsi permasalahan dan dampaknya.
  \item \textbf{Permasalahan Ketiga:} Deskripsi permasalahan dan dampaknya.
\end{enumerate}

\begin{figure}[H]
\centering
% \includegraphics[width=0.85\textwidth]{figures/diagram-masalah}
\fbox{\parbox{12cm}{\centering\vspace{2cm}[Diagram Identifikasi Masalah Bisnis]\vspace{2cm}}}
\caption{Diagram Identifikasi Masalah Bisnis}
\label{fig:diagram-masalah}
\end{figure}

\section{Solusi dan Rencana Implementasi}
\label{sec:solusi-implementasi}

\subsection{Alternatif Solusi}
\label{subsec:alternatif-solusi}

Berdasarkan analisis permasalahan yang telah dilakukan, terdapat beberapa alternatif
solusi yang dapat dipertimbangkan:

\begin{enumerate}
  \item \textbf{Alternatif 1:} Deskripsi alternatif solusi pertama beserta kelebihan
        dan kekurangannya.
  \item \textbf{Alternatif 2:} Deskripsi alternatif solusi kedua beserta kelebihan
        dan kekurangannya.
\end{enumerate}

\subsection{Rekomendasi Solusi}
\label{subsec:rekomendasi}

Berdasarkan evaluasi terhadap alternatif-alternatif solusi di atas, solusi yang
direkomendasikan adalah [nama solusi] karena \ldots

\subsection{Rencana Implementasi}
\label{subsec:rencana-implementasi}

Rencana implementasi solusi yang direkomendasikan adalah sebagai berikut:

\begin{table}[H]
\caption{Rencana Implementasi Solusi}
\label{tab:rencana-implementasi}
\begin{isitable}
\begin{tabular}{p{0.5cm} p{5cm} p{3cm} p{3cm} p{3cm}}
\toprule
\textbf{No} & \textbf{Kegiatan} & \textbf{Penanggung Jawab} & \textbf{Waktu} & \textbf{Indikator Keberhasilan} \\
\midrule
1 & Kegiatan pertama & Pihak penanggung jawab & Bulan 1-2 & Indikator \\
2 & Kegiatan kedua & Pihak penanggung jawab & Bulan 3-4 & Indikator \\
3 & Kegiatan ketiga & Pihak penanggung jawab & Bulan 5-6 & Indikator \\
\bottomrule
\end{tabular}
\end{isitable}
\end{table}

\section{Hasil Implementasi dan Evaluasi}
\label{sec:hasil-evaluasi}

Uraikan hasil implementasi solusi yang telah dilakukan (jika sudah diimplementasikan)
atau proyeksi dampak dari implementasi solusi yang direkomendasikan.

\begin{enumerate}
  \item \textbf{Dampak jangka pendek:} \ldots
  \item \textbf{Dampak jangka menengah:} \ldots
  \item \textbf{Dampak jangka panjang:} \ldots
\end{enumerate}
